\subsection{Matrix decomposition}\label{sec:matrix_decomp}

We start this section by defining the set of paths that will be relevant to 
bound the norm of $\|\bar{B}^{\ell}\|$.
We closely follow \cite{bordenave2015}.
%For this section, we assume $G$ is 
%$\ell$-tangle-free, which will hold with high probability.
\begin{definition}
Define $\Gamma^{\ell}_{ef}$ to be the set of all 
non-backtracking paths of $2\ell+1$ half edges, 
starting at $e$ and ending at $f$. 
A path in this set will be denoted by
$\gamma= (e_1,e_2,\dots,e_{2\ell},e_{2\ell+1})$,
where $e_1 = e$ and $e_{2\ell + 1} = f$. 
The non-backtracking property means that,
for all $1\leq i\leq \ell$, $e_{2i}$ and $e_{2i+1}$ share the same vertex 
but $e_{2i}\neq e_{2i+1}$.
Similarly, let $\Gamma^\ell = \bigcup_{e,f} \Gamma^\ell_{ef}$.
%this is to ensure we end on the right side. We use the transpose to concatenate!
\end{definition}

Each path in $\Gamma^\ell_{ef}$ uses $2 \ell + 1$ half edges,
corresponding to $\ell + 1$ edges in the graph.
To be clear, the above definition counts all possible
non-backtracking sequences of half edges.
These are different than the usual non-backtracking paths 
{\em and do not necessarily exist in the graph}.
Some of these paths might backtrack along a duplicate
edge which utilizes a different half edge.


We now have
%%%%%%%%%%%%%%%%%%%%%%%%%%%%%%%%%%%%%%%%%%%%%%%%%%%%%%%%%%%%%%%%%%%%%%%%
\begin{comment}
non-backtracking paths
of length $2\ell+1$, starting at half edge $e$ and ending at $f$.  
For a path $\gamma \in \Gamma^{\ell}_{ef}$,
we write
$\gamma = (e_1, e_2, \ldots, e_{\ell+1})$ 
where $e_i \in \vec{E}$ for all $i$,
$e_1 = e$ and $e_{\ell+1} = f$.
Similarly, define
$F^{\ell}_{ef} \subset \Gamma^\ell_{ef}$ 
be the set of all non-backtracking, tangle-free paths
in $G$
of length $\ell+1$, starting at oriented edge $e$ and ending at $f$. 
\end{comment}
%%%%%%%%%%%%%%%%%%%%%%%%%%%%%%%%%%%%%%%%%%%%%%%%%%%%%%%%%%%%%%%%%%%%%%%%%%
\begin{equation}
\label{eq:B_in_terms_of_M}
(B^{\ell})_{ef} = 
\sum_{\gamma \in \Gamma^\ell_{ef}}
\prod_{t=1}^{\ell} B_{e_{2t-1} e_{2t+1}}
= 
\sum_{\gamma \in \Gamma^\ell_{ef}}
\prod_{t=1}^{\ell} M_{e_{2t-1} e_{2t}}N_{e_{2t}e_{2t+1}}
=
\sum_{\gamma \in \Gamma^\ell_{ef}}
\prod_{t=1}^{\ell} M_{e_{2t-1} e_{2t}}
\end{equation}
%where we note the last equality requires $G$ to be $\ell$-tangle-free.
where we used Eqn.~\ref{eq:B=MN}
and the fact that $\Gamma^{\ell}_{ef}$ is non-backtracking, 
so
$
N_{e_{2t}e_{2t+1}}=1
$.

Recall that we will use
Eqn.~\eqref{eq:upper_bound_lambda2} and
Lemma \ref{lem:subspace} to bound $\lambda_2$.
Denote by $\bar{B}$ the matrix with entries equal to
$\bar{B} = B - S$,
where
\[
 S = \frac{1}{|E|}
 \left(
 \begin{array}{cc} 
 0 & (d_2 - 1) \1 \1^* \\
 (d_1 - 1) \1 \1^* & 0 
 \end{array}
 \right ).
\]
Note that $\bar{B}$ is an \textit{almost} centered version of $B$,
and $\Ker ( S ) = \Ker(T) = \mathrm{span} (\onealpha, \onemalpha ) $, 
where $T$ is the matrix from Lemma \ref{lem:subspace}.
To apply the lemma, we wish to get an expression like Eqn.~\eqref{eq:B_in_terms_of_M} for $\bar{B}^{\ell}$. To do so, we write:
\[
\bar{B} = B - S = M N - S = (M - S') N
\]
the above matrix equation in the unknown $S'$ can be solved by simple manipulations. We get
\[
S'=\frac{1}{|E|} \left(
 \begin{array}{cc} 
 0 & \1 \1^* \\
  \1 \1^* & 0 
 \end{array}
 \right ).
\]
Using again that $N$ is identically one over the elements of the set 
$\Gamma^{\ell}_{ef}$, 
we find a similar formula to Eqn.~\eqref{eq:B_in_terms_of_M}:
\begin{align}\label{eq:barB_in_terms_of_M}
(\bar{B}^\ell)_{ef}
=
\sum_{\gamma \in \Gamma^\ell_{ef}} 
\prod_{t=1}^{\ell}
\left( B - S \right)_{e_{2t-1} e_{2t+1}}
=
\sum_{\gamma \in \Gamma^\ell_{ef}} 
\prod_{t=1}^{\ell}
\bar{M}_{e_{2t-1} e_{2t}},
\end{align}
where $\bar{M}=M-S'$. 


The following telescoping sum formula is a simple algebraic manipulation 
and appears in \citet*{massoulie2013} 
and \citet*{bordenave2015a}: 
\begin{equation*}
\prod_{s=1}^\ell x_s = \prod_{s=1}^\ell y_s + 
\sum_{j=1}^\ell \prod_{s=1}^{j-1} y_s (x_j - y_j) \prod_{t=j+1}^\ell x_t .
\end{equation*}
Using this, 
with $x_s = B_{e_{2s-1} e_{2s+1}}$ 
and $y_s = \bar{B}_{e_{2s-1} e_{2s+1}}$,
we obtain the following relation:
\begin{align}
\label{B_dec1}
( B^{\ell} )_{ef} 
= 
( \bar{B}^{\ell} )_{ef} 
+ 
\sum_{\gamma\in \Gamma^{\ell}_{ef}} 
\sum_{j=1}^{\ell} 
\prod_{s=1}^{j-1}\
\bar{B}_{e_{2s-1} e_{2s+1}}
S_{e_{2j-1} e_{2j+1}}
\prod_{t=j+1}^{\ell} B_{e_{2t-1} e_{2t+1}}.
\end{align}
This decomposition breaks the elements in $\Gamma^{\ell}_{ef}$ 
into two subpaths, also non-backtracking, 
of length $j$ and $\ell-j$, respectively. 

\begin{definition}
Let $F_{ef}^\ell \subset \Gamma^\ell_{ef}$ 
denote the subset of paths
which are tangle-free,
with $F^\ell = \bigcup_{e,f} F^\ell_{ef}$. 
\end{definition}

We will take the parameter $\ell$ to be small enough so that
the path $\gamma$ is tangle-free with high probability.
Thus the sums in Eqns.~\eqref{eq:B_in_terms_of_M}
or \eqref{eq:barB_in_terms_of_M}
need only be over the paths 
$\gamma \in F_{ef}^\ell$.
However,
to recover the matrices $B$ and $\bar{B}$ by rearranging 
Eqn.~\eqref{B_dec1},
we need to also count those tangle-free 
subpaths that arise from splitting tangled paths.
While breaking a tangle-free path will necessarily
give us two new tangle-free subpaths, the converse is not always true. 
This extra term generates a remainder that we define now.

\begin{definition}
\label{def:tangled_paths}
Let $T^{\ell,j}_{ef}$ 
be the set of non-backtracking paths containing $2\ell+1$ half edges,
starting at $e$ and ending at $f$,
such that overall the path is tangled
but the first $2j-1$, middle three, and last $2(\ell - j)+1$ 
half edges form tangle-free subpaths:
$\gamma = (e_1, \ldots, e_{2\ell + 1}) \in T^{\ell,j}$
if and only if
$\gamma' = (e_1, \ldots, e_{2j-1}) \in F^{j-1}$,
$\gamma'' = (e_{2j-1}, e_{2j}, e_{2j+1}) \in F^{1}$,
and $\gamma''' = (e_{2j+1}, \ldots, e_{2\ell + 1}) \in F^{\ell - j}$.
Set $T^{\ell,j} = \bigcup_{e,f} T^{\ell,j}_{ef}$.
\end{definition}
Set the remainder
\begin{align}
    \label{error_term}
    R^{\ell,j}_{ef} 
    &= 
    \sum_{\gamma\in T^{\ell,j}_{ef}} 
    \sum_{j=1}^{\ell}
    \prod_{s=1}^{j-1}
    \bar{B}_{e_{2s-1} e_{2s+1}}
    S_{e_{2j-1} e_{2j+1}}
    \prod_{t=j+1}^{\ell} B_{e_{2t-1} e_{2t+1}}
    \nonumber \\
    &=
    \sum_{\gamma\in T^{\ell,j}_{ef}} 
    \sum_{j=1}^{\ell} 
    \prod_{s=1}^{j-1}\
    \bar{M}_{e_{2s-1} e_{2s}}
    S_{e_{2j-1} e_{2j+1}}
    \prod_{t=j+1}^{\ell} 
    M_{e_{2t-1} e_{2t}}.
\end{align}
Since the paths are non-backtracking, 
the $N$ terms are all unity.

Adding and subtracting 
$\sum_{j=1}^{\ell} R^{\ell,j}_{ef}$ 
to Eqn.~\eqref{B_dec1}
and rearranging the sums, we obtain
\begin{align}
\label{B_dec2}
B^{(\ell)} = 
\bar{B}^{(\ell)} + 
\sum_{j=1}^{\ell} 
\bar{B}^{(j)} S B^{(\ell-j)}
- 
\sum_{j=1}^{\ell} R^{\ell,j},
\end{align}
where the matrices $B^{(\ell)}$ and $\bar{B}^{(\ell)}$
are tangle-free versions of $B^\ell$ and $\bar{B}^\ell$,
i.e.\ element $ef$ in both matrices
only counts paths $\gamma \in F_{ef}^\ell$.
Multiplying Eqn.~\eqref{B_dec2} on the right by 
$x \in \Ker (T)$
and using that 
$B^{(\ell - j)} x$
is also within $\Ker(S)$,
since it is just the space spanned by the leading eigenvectors,
we find that the middle term is identically zero.
Thus for $x \in \Ker(T)$,
\begin{equation}\label{decomposition}
\| B^{(\ell)}x \| 
\leq 
\| \bar{B}^{(\ell)} x \|+
\left\| 
%\|
\sum_{j=1}^{\ell}R^{\ell,j} x 
%\|
\right\|
\leq 
\| \bar{B}^{(\ell)} \|+
\sum_{j=1}^{\ell} \| R^{\ell,j} \| .
\end{equation}
